% !PRINTED 59
\weberpage{59}

%<<CONT
unity, then $e$ is by (1) a divisor of $2^{\lambda-2}$, hence a power of 2, and if $e < 2^{\lambda-2}$, then we should have to have

\[
5^{2^{\lambda-3}} \equiv 1 \pmod{2^\lambda}
\]

But this is not possible, for, for every $\lambda$ greater than 2, the formula

\[
5^{2^{\lambda-3}} = 1 + 2^{\lambda-1} h\ ^{1)}
\]

holds with odd $h$ --- a formula which, just like the formula (1), can be proved by complete induction. Hence, if $\lambda \geqq 3$,
the $2^{\lambda-2}$ powers

\weq{2}{1, 5, 5^2 \ldots 5^{2^{\lambda-2}-1}}

are all different from one another with respect to the modulus $2^\lambda$. Now a relation of the form $5^h \equiv -5^\varkappa$ is impossible for the modulus 4, and therefore all the more for every higher power of 2, and consequently the quantities

\weq{3}{-1, -5, -5^2 \ldots -5^{2^{\lambda-2}-1}}

are, with respect to the modulus $2^\lambda$, different both among themselves and from the quantities (2); and since their total number amounts to $2^{\lambda-1}$, every odd number is congruent, with respect to the modulus $2^\lambda$, to one and only one of the quantities (2), (3).

If, then, we set

\weq{4}{a = 2, \quad b = 2^{\lambda-2},}

we obtain the following theorem for $\lambda \geqq 3$:

\noindent\textbf{3.} \gesperrt{For every odd number} $n$ \gesperrt{one can assign, completely determined with respect to the pair of moduli} $a, b$\gesperrt{, a pair of numbers} $\alpha, \beta$\gesperrt{, so that}

\[
n \equiv (-1)^\alpha 5^\beta \pmod{2^\lambda}
\]

\gesperrt{holds.}

The case $\lambda = 2$ can be subsumed under this, because then $b = 1$ and $\beta = 0$ can be set. In order to include the case $\lambda = 1$ under it as well, which however offers no particular interest, one would have to set $a = 1, b = 1$.

\noindent\rule{4cm}{0.4pt}

\noindent $^{1)}$ If one wished to take the number 3 as base in place of the number 5, this formula would not yet be valid for $\lambda = 3$, though it would be for every larger $\lambda$, and hence 3 could serve as base just as well as 5. The reason for preferring the base 5 lies in the fact that all powers of this base are $\equiv 1 \pmod{4}$.
